\documentclass{article}
\usepackage{amsmath, amssymb, amsthm}
\usepackage{hyperref}
\usepackage{geometry}
\geometry{margin=1in}

\title{Erd\H{o}s Problem \#170}
\date{Last updated: 2025-08-31}
\author{Source: erdosproblems.com}

\begin{document}
\maketitle

\noindent\textbf{Status:} OPEN \quad \textbf{No prize}

\noindent\textbf{Tags:} additive combinatorics

\medskip
\noindent\textbf{Problem Statement:}

Let $F(N)$ be the smallest possible size of $A\subset \{0,1,\ldots,N\}$ such that $\{0,1,\ldots,N\}\subset A-A$. Find the value of\[\lim_{N\to \infty}\frac{F(N)}{N^{1/2}}.\]

\bigskip
\noindent\textbf{Remarks:}

The Sparse Ruler problem. R\'{e}dei asked whether this limit exists, which was proved by Erd\H{o}s and G\'{a}l \cite{ErGa48}. Bounds on the limit were improved by Leech \cite{Le56}. The limit is known to be in the interval $[1.56,\sqrt{3}]$. The lower bound is due to Leech \cite{Le56}, the upper bound is due to Wichmann \cite{Wi63}. Computational evidence by Pegg \cite{Pe20} suggests that the upper bound is the truth. A similar question can be asked without the restriction $A\subset \{0,1,\ldots,N\}$.

\bigskip
\noindent\textbf{References:}

[ErGa48] Erd\H{o}s, P. and G\'{a}l, I., On the representation of $1,2,\ldots,N$ by differences. Nederl. Akad. Wetensch., Proc. (1948), 1155-1158.
 
 [Le56] Leech, J., On the representation of $1,2,\ldots,n$ by differences. J. London Math. Soc. (1956), 160-169.
 
 [Pe20] Pegg, E., Hitting All the Marks: Exploring New Bounds for Sparse Rulers and a Wolfram Language Proof. https://blog.wolfram.com/2020/02/12/hitting-all-the-marks-exploring-new-bounds-for-sparse-rulers-and-a-wolfram-language-proof/ (2020).
 
 [Wi63] Wichmann, B., A note on restricted difference bases. J. London Math. Soc. (1963), 465-466.

\bigskip
\noindent\small{Source: \url{https://www.erdosproblems.com/170}}
\end{document}
